Let $A=1$ mean that the gate accepts the initial concept. After a shared
prefix of cost $C_p$, let $C_e$ and $C_s$ be the remaining costs of the
early and search branches, including their final segmentation calls.
The gated and staged always-search costs satisfy
\begin{equation}
  \begin{aligned}
  C_g&=C_p+A C_e+(1-A)C_s,\\
  C_{\mathrm{st}}&=C_p+C_s.
  \end{aligned}
  \label{eq:gate-cost}
\end{equation}
Writing $D=C_s-C_e$ and denoting the corresponding final-mask IoUs by
$J_e$ and $J_s$ gives
\begin{equation}
  \begin{aligned}
  \mathbb{E}[C_g-C_{\mathrm{st}}]&=-\mathbb{E}[A D],\\
  \mathbb{E}[J_g-J_{\mathrm{st}}]&=\mathbb{E}[A(J_e-J_s)].
  \end{aligned}
  \label{eq:gate-decomposition}
\end{equation}
The first identity measures computation avoided on accepted inputs;
the second measures the accuracy change on those same inputs. If
$H=C_{\mathrm{st}}-C_{\mathrm{full}}$ is the staging overhead relative to
original CG-ICS, the gate is faster in expectation precisely when
$\mathbb{E}[A D]>\mathbb{E}[H]$. Consequently, an acceptance rate alone
does not establish a speedup: the saved branch costs must exceed the
probe overhead. The staged always-search and random-routing controls
separate this overhead from the value of the agreement signal.

